
This study set out to investigate ideological biases in large language models by analyzing agreement ratios and the semantic similarity of justifications across topics, roles, and languages. The findings highlight several key tendencies.

Across both GPT-4o-mini and LLaMA~3.1-70B, certain topics emerged as more stable than others. For instance, \emph{Multi-party system} and, often, \emph{Adopt atheism} consistently showed high similarity across conditions, suggesting a stronger shared framing or more globally uniform reasoning. In contrast, topics such as \emph{No economic sanctions}, \emph{Compulsory voting}, and \emph{Adopt libertarianism} displayed lower similarity, indicating that the models diverged more strongly in their reasoning whenever a persona prompt was applied. These results underline that ideological variation is highly topic-sensitive rather than evenly distributed.

A particularly important observation is that persona prompts influence outputs more than language changes. The \emph{Role vs. Assistant} comparisons (RvA) consistently exhibited greater divergence than either \emph{Native vs. Assistant} (NvA) or \emph{Native .vs Role} (NvR). This shows that assigning a national identity or role has a stronger effect on model reasoning than merely switching the input language. While multilinguality shapes nuances of phrasing, persona conditioning actively reshapes stance and justification. This distinction is central to understanding how users might inadvertently steer models when framing prompts.

When comparing across countries, a relatively stable structure could be observed. The USA remained among the highest and most consistent across topics, Western European countries such as France, Italy, Germany, and Russia occupied a mid–high range, while China, Israel, Iran, and Ukraine often appeared lower and more topic-sensitive. This suggests that models implicitly reproduce certain cultural and ideological framings embedded in their training data, which cluster countries into identifiable patterns.

Comparing the two models further emphasized the role of alignment and fine-tuning. Both GPT and LLaMA displayed the same broad topic sensitivities, but LLaMA exhibited stronger persona-driven divergence. GPT responses were more uniform, likely reflecting heavier reinforcement learning and alignment procedures. This indicates that ideological bias in LLMs is not solely determined by pretraining corpora but is also shaped by post-training adjustments.

Taken together, these findings emphasize the importance of systematic evaluation of ideological bias in LLMs. Agreement ratios revealed how stance distribution varies across roles and languages, while embedding similarities showed that divergences persist even when surface-level agreement remains. The evidence that persona conditioning exerts stronger influence than language variation has direct implications for global deployment: models may appear linguistically fluent yet deliver role-dependent reasoning that diverges ideologically. In sensitive domains such as journalism, policy-making, or education, such hidden variation reinforces the need for transparency, careful benchmarking, and continuous monitoring of ideological bias.
